\documentclass[10pt,a4paper]{amsart}
\usepackage[margin=19mm]{geometry}
\usepackage{amsmath,amssymb,mathtools}
\usepackage[scr=rsfs,cal=euler]{mathalfa}
\usepackage{xfrac}
\usepackage[unicode,hidelinks,bookmarks=false]{hyperref}
\newcommand{\ie}{i.e.\ }
\newcommand{\cf}{cf.\ }
\newcommand{\resp}{respectively\ }
\newcommand{\Subsection}{\subsection}
\providecommand{\nobreakdash}{\nobreak\hbox{-}}
\setlength{\emergencystretch}{2em}
\multlinegap0pt
\usepackage{xcolor}
\usepackage{amssymb}
\usepackage{algorithm}\usepackage{algpseudocode}
\newtheorem{lemma}{Lemma}
\title{The Computational Complexity of Counting Linear Regions in ReLU Neural Networks}
\author{Moritz Stargalla, Christoph Hertrich, Daniel Reichman}
\date{Source: arXiv:2505.16716v3 (CC BY 4.0)}
\begin{document}
\maketitle

\noindent\emph{Source and licence.} The Computational Complexity of Counting Linear Regions in ReLU Neural Networks. Version arXiv:2505.16716v3 (CC BY 4.0), distributed by arXiv under the Creative Commons Attribution 4.0 International licence, http://creativecommons.org/licenses/by/4.0/, as recorded in the arXiv licence metadata for this exact version and shown on the abstract page https://arxiv.org/abs/2505.16716v3. Full licence notice: \texttt{licenses/arxiv-2505.16716-CC-BY-4.0.txt}.\par\smallskip

\noindent\textit{Editorial scope.} Counted unit: the article's lemma lem:coeffs\_are\_polynomial (source0.tex lines 604--611). The article's complete original proof of it is delivered with it (source0.tex lines 612--640, one \texttt{proof} environment of 29 lines). No mathematics is written, repaired or re-derived: the statement and the proof are the article's own lines, in the article's own notation, and no retained line is substituted or re-flowed.  No further source-local context is needed: every cross-reference of every retained line resolves inside the two delivered spans.  Nothing was omitted from any retained span, so the package carries no omission record.  No retained line contains a citation, so the wrapper carries no bibliography and no bibliography entry.  No retained line contains a figure environment, an image-inclusion command or drawing code, so no image is delivered and the package declares no asset dependency.  Editorial formatting only, in addition: an AMS \texttt{amsart} preamble that restates the article's own declarations and macros used by the retained lines, namely the environments \texttt{lemma} on separate counters, together with the standard packages those lines need.  The retained environments print in this document as Lemma 1 in order of appearance, whereas the article prints the same items with the numbering of its own full document.  The complete native source of the article as distributed in the arXiv e-print for this exact version is bundled unchanged as \texttt{sources/178585/source0.tex}.
\begin{lemma}
\label{lem:coeffs_are_polynomial}
    Given a ReLU network $N$ and an activation pattern $a\in \{0,1\}^{s(N)}$ corresponding to a proper activation region. Let $A_{\max}\in \mathbb{Z}$ be the maximum absolute value of any numerator or denominator in an entry of the matrices $A^{(1)},\dots, A^{(d+1)}$ and biases $b^{(1)},\dots, b^{(d+1)}$ and let $n_{\max}=\max\{n_0,\dots,n_{d+1}\}$.
    Then, the encoding size of every coefficient of the affine function $f_N^a:\mathbb{R}^n\to \mathbb{R}$, $f_N^a(x)=\sum_{i=1}^n a_ix + b$ is bounded by
    \[
    36d^2n_{\max}^2\langle A_{\max}\rangle.
    \]
\end{lemma}

\begin{proof}
    Given an activation pattern $a\in \{0,1\}^{s(N)}$ of $N$, we modify the matrices $A^{(1)},\dots, A^{(d+1)}$ and biases $b^{(1)},\dots, b^{(d+1)}$ by replacing the columns of matrices and bias entries corresponding to inactive neurons with 0 entries.
    Let $A^{(1),a},\dots, A^{(d+1),a}$ and biases $b^{(1),a},\dots, b^{(d+1), a}$ denote the modified matrices and biases.
    Then, a simple calculation shows that for all $x\in \mathbb{R}^n$, we have
    \[
    f_N^a(x) = A^{(d+1),a}\cdot \dots \cdot A^{(1),a}x + b^{(d+1), a}+\sum_{i=0}^{d-1}A^{(d+1),a}\cdot \dots \cdot A^{(d+1-i),a}b^{(d-i), a}.
    \]
    To bound the encoding size of all occurring coefficients, we separately give a bound on the absolute value of any occurring denominator and numerator.
    For this, we turn the $d+1$ rational matrices and biases into integral matrices and biases by bringing all fractional entries to a common denominator.
    The value of the common denominator is bounded by $A_{\max}^{(d+1)(n_{\max}+1)^2}$, since there are fewer than $(d+1)(n_{\max}+1)^2$ entries in the rational matrices and biases.

    To bound the maximum absolute numerator value, we now consider the $(d+1)$ integral matrices and biases that arise by multiplying the fractional matrices by the common denominator.
    The maximum absolute value of an entry in one of the integral matrices or biases is bounded by $A_{\max}^{(d+1)(n_{\max}+1)^2}$.
    A simple calculation shows that the maximum absolute entry that can be obtained in the product of the $(d+1)$ integral matrices is
    \[
    \left(A_{\max}^{(d+1)(n_{\max}+1)^2}\right)^{d+1}\prod_{i=1}^{d}n_i
    \leq n_{\max}^d \cdot A_{\max}^{(d+1)^2(n_{\max}+1)^2}.
    \]
    The constant of $f_N^a$ is equal to $b^{(d+1), a}+\sum_{i=0}^{d-1}A^{(d+1),a}\cdot \dots \cdot A^{(d+1-i),a}b^{(d-i), a}$ and can thus be bounded by $(d+1) n_{\max}^d  A_{\max}^{(d+1)^2(n_{\max}+1)^2}$.
    
    Since the maximum encoding size of an element of a set of integers is obtained by the integer having the maximum absolute value, it follows that the encoding size of a coefficient in $f_N^a$ is at most
    \begin{align*}
    &1+\langle (d+1) n_{\max}^d  A_{\max}^{(d+1)^2(n_{\max}+1)^2} \rangle + \langle A_{\max}^{(d+1)(n_{\max}+1)^2} \rangle\\
    \leq\;&1+\langle d + 1 \rangle+d\langle n_{\max}\rangle+(d+1)^2(n_{\max}+1)^2\langle A_{\max}\rangle+(d+1)(n_{\max}+1)^2\langle A_{\max}\rangle\\
    \leq\;&1+\langle d + 1 \rangle+d\langle n_{\max}\rangle+2(d+1)^2(n_{\max}+1)^2\langle A_{\max}\rangle\\
    \leq\;&1+\langle d + 1 \rangle+d\langle n_{\max}\rangle+32d^2n_{\max}^2\langle A_{\max}\rangle\\
    \leq \; & 36d^2n_{\max}^2\langle A_{\max}\rangle.
    \end{align*}
\end{proof}

\end{document}
